\chapter[\emj{📈}~DP: Longest Increasing Subsequence (LIS)]{\emj{📈}~DP: Longest Increasing Subsequence (LIS)}

\begin{dsaquees}

De una lista de números 📊, encontrar la subsecuencia CRECIENTE más larga:
los números van subiendo (no importa si te saltas algunos, pero deben ir
en orden y de menor a mayor).

Ejemplo: en [10, 9, 2, 5, 3, 7, 101, 18] una subida larga es [2, 3, 7,
18] o [2, 5, 7, 101], ambas de largo 4. Hay dos formas de resolverlo: una
simple $O(n^2)$ y una astuta $O(n \log n)$ con búsqueda binaria.

\end{dsaquees}

\begin{dsanino}

Imagina que anotas tu estatura 📏 cada año, pero en desorden. Quieres
encontrar la racha más larga en la que fuiste creciendo, aunque tengas
que ignorar algunos años.

La forma fácil: para cada número te preguntas "¿cuál es la subida más
larga que TERMINA justo aquí?". Miras todos los números anteriores más
pequeños y te quedas con la mejor racha, más uno.

La forma lista: vas armando una "fila de paciencia". Cada número nuevo
reemplaza al primer número de la fila que sea igual o más grande que él
(así la fila queda lo más chica posible). Al final, el LARGO de esa fila
es la respuesta. Es como un solitario de cartas donde apilas de forma
ordenada.

\end{dsanino}

\begin{dsaotraforma}

El LIS busca la subsecuencia estrictamente creciente más larga de un
arreglo (los elementos conservan su orden original, no tienen que ser
contiguos).

\textbf{Enfoque DP $O(n^2)$}: \verb|dp[i]| = longitud del LIS que termina
en el índice $i$. Para cada $i$, mira todos los $j < i$ con
\verb|nums[j] < nums[i]| y toma \verb|dp[i] = max(dp[i], dp[j] + 1)|.

\textbf{Enfoque $O(n \log n)$}: mantén un arreglo \verb|tails| donde
\verb|tails[k]| es el menor valor con el que puede terminar una
subsecuencia creciente de longitud $k+1$. Por cada número, usa búsqueda
binaria (\verb|lower_bound|) para colocarlo; el tamaño final de
\verb|tails| es la respuesta. Ojo: \verb|tails| NO es el LIS real, solo
su longitud.

\end{dsaotraforma}

\begin{dsadiagrama}
nums = [10, 9, 2, 5, 3, 7, 101, 18]

Método O(n log n) construyendo 'tails' (lower_bound reemplaza):

  10 -> tails=[10]
   9 -> reemplaza 10   -> [9]
   2 -> reemplaza 9    -> [2]
   5 -> agrega         -> [2,5]
   3 -> reemplaza 5    -> [2,3]
   7 -> agrega         -> [2,3,7]
 101 -> agrega         -> [2,3,7,101]
  18 -> reemplaza 101  -> [2,3,7,18]

Longitud de tails = 4  => LIS = 4
(tails = [2,3,7,18] NO es la subsecuencia real, pero su LARGO sí es 4)
\end{dsadiagrama}

\begin{lstlisting}
CÓMO FUNCIONA (O(n log n))
- tails vacío. Por cada x:
  * Busca con lower_bound la primera posición con valor >= x.
  * Si no existe (x es mayor que todos) -> agrega x al final.
  * Si existe -> reemplaza ese valor por x (mantiene tails "apretado").
- La longitud de tails es el LIS.

PSEUDOCÓDIGO (O(n log n))
──────────────────────────────────────────────────────────────

función lis(nums):
    tails = []
    PARA cada x en nums:
        pos = lower_bound(tails, x)     // primer >= x
        SI pos == fin(tails): tails.push(x)
        SINO: tails[pos] = x
    return longitud(tails)

CÓDIGO C++ (DP O(n^2), fácil de entender)
──────────────────────────────────────────────────────────────

#include <vector>
#include <algorithm>
using namespace std;

int lisDP(vector<int>& nums) {
    int n = nums.size();
    vector<int> dp(n, 1);                 // cada elemento es un LIS de 1
    int best = n ? 1 : 0;
    for (int i = 1; i < n; i++)
        for (int j = 0; j < i; j++)
            if (nums[j] < nums[i]) {
                dp[i] = max(dp[i], dp[j] + 1);
                best = max(best, dp[i]);
            }
    return best;
}

CÓDIGO C++ (O(n log n) con búsqueda binaria)
──────────────────────────────────────────────────────────────

int lisFast(vector<int>& nums) {
    vector<int> tails;
    for (int x : nums) {
        auto it = lower_bound(tails.begin(), tails.end(), x);
        if (it == tails.end()) tails.push_back(x);  // extiende
        else *it = x;                               // reemplaza
    }
    return tails.size();
}

COMPLEJIDAD (Big-O)
──────────────────────────────────────────────────────────────

  Enfoque       │ Tiempo       │ Espacio
  ──────────────┼──────────────┼──────────
  DP clásico    │ O(n^2)       │ O(n)
  Binary search │ O(n log n)   │ O(n)

* Para LIS no decreciente (permite iguales), usa upper_bound.
\end{lstlisting}

\begin{dsaentrevistas}

✅ "Longest Increasing Subsequence" (LC 300): si te piden solo la
   LONGITUD, ve directo al $O(n \log n)$. Si piden reconstruir la
   subsecuencia real, el DP $O(n^2)$ con arreglo de padres es más claro.

💡 Estricto vs no decreciente: \verb|lower_bound| da LIS estrictamente
   creciente; \verb|upper_bound| permite elementos iguales (no
   decreciente). Un detalle que cambia la respuesta.

⚠️ tails no es la subsecuencia: mucha gente cree que \verb|tails| ES el
   LIS. Solo su LONGITUD es correcta; los valores pueden no formar una
   subsecuencia válida.

🔑 Disfraces del LIS: "Russian Doll Envelopes", "Number of LIS",
   "Longest Chain of Pairs" son LIS con un ordenamiento previo. Reconocer
   el patrón vale oro.

\end{dsaentrevistas}

\begin{dsaresumen}

• LIS: subsecuencia creciente más larga (respeta orden, permite huecos).
• DP O(n²): dp[i] = mejor LIS que termina en i; mira los j < i menores.
• O(n log n): arreglo tails + \verb|lower_bound|; su longitud es el LIS.
• tails da la LONGITUD correcta, no la subsecuencia real.
• Estricto -> \verb|lower_bound|; no decreciente -> \verb|upper_bound|.
• Muchos problemas se reducen a LIS tras ordenar.
════════════════════════════════════════════════════════════════

\end{dsaresumen}
